% ============================================================================
\section{Introduction}
\label{sec:intro}
% ============================================================================

Multi-teacher on-policy distillation (MOPD) provides a practical way to
combine capabilities in one language model. The student generates responses,
and each response is scored by the teacher assigned to its task
\citep{ma2026mopdmultiteacheronpolicydistillation}. This allows specialist
teachers to be developed separately and their feedback to be combined during
student training. The same approach can add a new capability while using an
earlier model as a teacher to preserve existing ones
\citep{liu2026instella}. The intended outcome is a student that acquires
useful capabilities across tasks. A decrease in average teacher loss alone
does not establish that outcome.

Three questions must be distinguished. Can this student acquire each
capability through separate distillation? Can one shared model attain
those capability levels together? Can training reach them with the
available computation? Open-MOPD uses separately distilled students as
references and identifies capability imbalance in joint training
\citep{gao2026openmopddiagnosingfixingcapability}. These references establish
what the student learns separately; a gap in joint training can still
reflect task exposure, interacting updates, or other optimization effects.

We study a local question about the training process:
\begin{quote}
Given capabilities that a student can acquire separately, when does a
shared MOPD update improve task objectives together, and how much
computation is needed to make that progress reliable?
\end{quote}
A task can receive a harmful update for two distinct reasons. The weighted
mean gradient after optimizer scaling may oppose its loss gradient, or a
useful mean direction may be estimated too imprecisely. Additional samples
address the second cause.
Under fixed task weights and optimizer scaling, they cannot change the
expected first-order effect of the mean direction.

Existing work motivates testing both causes. CaMOPD studies counteracting
recovery and preservation gradients \citep{chen2026camopd}; D$^3$-MOPD
allocates samples using the remaining teacher gap and descent velocity
\citep{sun2026d3mopd}. Slow alignment has also been observed when the
training states are held fixed \citep{fu2026rethinkingopd2}. These findings
do not determine whether sampling uncertainty limits a particular MOPD
run. Dense top-$k$ supervision already computes multiple token
contributions at each prefix, while prompts and student trajectories
continue to produce variable gradients
\citep{ma2026mopdmultiteacheronpolicydistillation}.

Our analysis measures each task's alignment with the combined mean update
and the uncertainty along its loss gradient. A classical one-sided
probability bound then relates sampling to the reliability of first-order
progress. The analysis concerns losses on fixed student-generated states;
we separately measure the effect of changing the states visited by the
student. Positive local progress is a diagnostic, not a requirement that
every task improve at every step of a successful training run.

Gradient-Preconditioned Adaptive Sampling (GPAS) provides an allocation
intervention for this analysis. It applies Neyman allocation to gradient
variance after optimizer scaling and preserves the specified task mean
weights. Its per-step rule minimizes an estimate of total update variance,
a conservative measure of uncertainty in task progress. We give an
analytical example in which the
variance-minimizing allocation has greater directional uncertainty than
another allocation. The experiment protocol tests when this distinction
affects actual updates and downstream capability.

The paper develops three components:
\begin{itemize}
    \item a local analysis separating mean task interactions from sampling
    uncertainty, with sufficient conditions for reliable joint descent;
    \item a characterization of what GPAS's total-variance criterion controls,
    including a counterexample to directional optimality;
    \item a four-domain experimental protocol testing these quantities
    against controlled updates, separately acquired capabilities, and
    end-to-end training efficiency.
\end{itemize}
\missing{Replace the protocol contribution with measured findings after
the diagnostic and training comparisons are complete.}
